(Notation correction.) The statement defines $F(n)$ using $\{1,\dots,N\}$ but then uses $n$ elsewhere; we interpret this as $A\subseteq\{1,\dots,n\}$.

For $n\in\mathbb N$, let $F(n)$ be the maximum possible size of a set $A\subseteq\{1,\dots,n\}$ with the property:
\[
\text{for all distinct }a,b,c,d\in A,\ \{a,b\}\ne\{c,d\}\implies ab\ne cd.
\]
Equivalently, the map $\{a,b\}\mapsto ab$ from unordered pairs of distinct elements of $A$ is injective.

Question 1: Does there exist a constant $c$ such that
\[
F(n)=\pi(n)+\bigl(c+o(1)\bigr)\,n^{3/4}(\log n)^{-3/2}\quad (n\to\infty)?
\]
The problem statement records Erd\H{o}s’s bounds: there exist constants $0<c_1\le c_2$ with
\[
\pi(n)+c_1 n^{3/4}(\log n)^{-3/2}\le F(n)\le \pi(n)+c_2 n^{3/4}(\log n)^{-3/2}.
\]

Question 2 (higher products): Fix an integer $r\ge 2$. If $A\subseteq\{1,\dots,n\}$ is such that for every two distinct increasing $r$-tuples $(a_1<\cdots<a_r)$ and $(b_1<\cdots<b_r)$ from $A$ we have
\[
 a_1\cdots a_r\ne b_1\cdots b_r,
\]
is it true that
\[
|A|\le \pi(n)+O\bigl(n^{(r+1)/(2r)}\bigr)?
\]
